\section{From intrinsic channel laws to the periodic table}
\label{sec:global}
A local contact law does not register different experiments in a common
plane. This section specifies finite intrinsic data that do so. The
registration consists of boundary distances and integer copy labels,
not positions, angles, curvatures, or known obstacle images.

\begin{definition}[Registered normal-channel network]\label{def:network}
Let $V$ index all the obstacle components in one cell. A finite directed
multigraph has an edge $e=(v,w,k_e)$ for each selected isolated clear
normal channel from a chosen copy of obstacle $v$ to the copy of obstacle
$w$ displaced by $Lk_e$, where $k_e\in\mathbb Z^2$ and $L$ is the
unknown lattice basis matrix. Loops and multiple edges are allowed.
The integer labels are part of the observation. The underlying graph
on $V$ is connected, and every vertex has an incident channel.

At each vertex distinguish one incident contact occurrence. For every
other incident occurrence supply its signed arclength distance from that
reference along the same boundary, modulo its perimeter. All distances
and all channel-law representatives use one coherent choice of boundary
orientation, considered modulo its single simultaneous reversal. At the
source endpoint a channel's transverse arclength has that boundary
orientation, and at the target it has the opposite orientation, as in
Section~\ref{sec:introduction}. Contact occurrences, not just edges,
are registered, including the two occurrences of a loop. No independent
sign reversal at a vertex or edge is allowed.

For a closed graph walk the signed sum of its integer edge labels is
its cycle label. The network has \emph{full lattice rank} if these labels
span $\R^2$.
\end{definition}

The orientation in this definition is intrinsic bookkeeping. Equality
of registered data means equality after one common sign choice, not
equality of independently symmetrized edge laws. Perimeters need not
be supplied separately: they are determined when a participating
analytic boundary image has been reconstructed. The registration must
be compatible with actual arclength on that image; the theorem assumes
the data come from physical tables.

\begin{theorem}[Registered periodic rigidity]\label{thm:network}
Let two periodic dispersing tables have connected closed real-analytic
obstacle boundaries. If their registered normal-channel networks have
full lattice rank and their intrinsic channel laws and registration agree,
then one Euclidean isometry takes the whole first labelled periodic table
and its marked lattice to the second. The lattice need not be supplied
numerically. The conclusion uses only the compressed local data in
\eqref{eq:arcdatum}--\eqref{eq:betadef}, their onsets, and the finite
registration of Definition~\ref{def:network}.
\end{theorem}
\begin{proof}
Choose matching representatives of the single orientation orbit. Fix
one vertex and one incident contact, and place that contact and its
outward normal in a numerical reference frame. Theorem~\ref{thm:intrinsic}(i)--(ii),
already proved without using this section, reconstructs its entire
analytic boundary image. Translation and rotation of the initial frame
are the only choices at this step. Arclength registration on this now
known closed curve locates every other incident contact uniquely modulo
the perimeter. In particular, it supplies each contact point and its
outward unit normal without measuring an ambient frame.

Choose a spanning tree rooted at this vertex. Assign integers
$\tau_v\in\mathbb Z^2$ along the tree so that each tree edge has
transformed label
\[
 k'_e=k_e+\tau_v-\tau_w=0.
\]
This only changes the chosen representative copy of each obstacle.
Traverse the tree. When a source contact is placed at $p_{v,e}$ with
outward normal $n_{v,e}$, the partner contact is at
\begin{equation}\label{eq:placepartner}
 p_{v,e}+g_e n_{v,e}.
\end{equation}
Its outward normal is $-n_{v,e}$. The recovered local graph and the
coherent arclength orientation place the partner germ uniquely there;
Lemma~\ref{lem:analyticidentity} places its entire boundary image.
Its registered contacts can then be located. If a tree edge is traversed
backwards, use its target occurrence and the reversed normal segment;
this is the same construction with the endpoints interchanged. Induction
therefore places all representative obstacles in one common plane.
For two tables with matching data, these placed images coincide.

For a remaining edge, both contact occurrences are now known on their
representative images. Its physical displacement must satisfy
\begin{equation}\label{eq:cyclelinear}
 d_e:=p_{v,e}+g_e n_{v,e}-p_{w,e}=L k'_e.
\end{equation}
The non-tree transformed labels generate the real span of all cycle
labels: adjoining a non-tree edge to the tree path gives precisely its
fundamental cycle, with label $k'_e$. Full rank supplies two edges
$e_1,e_2$ whose transformed labels are linearly independent. Hence
\begin{equation}\label{eq:latticerecover}
 L=[d_{e_1}\ d_{e_2}]
                 [k'_{e_1}\ k'_{e_2}]^{-1}.
\end{equation}
The integer matrix need not be unimodular. Its entries are known copy
labels, so invertibility over $\R$ suffices. Every other cycle gives
a consistency equation, not a further choice of periods. Thus the two
tables have the same marked lattice and the same obstacle images after
the initial common placement. Reversing the coherent orientation changes
this placement by one reflection. Restoring the arbitrary initial
translation and rotation proves the isometry assertion.
\end{proof}

This proves Theorem~\ref{thm:intrinsic}(iii). The graph hypothesis is a
finite visibility and registration condition, not a finite-horizon
assumption. It permits the two contacts of a channel to lie on distinct
translates of the same obstacle.

\begin{proposition}[Physical examples and the rank condition]
\label{prop:networkexamples}
There is a nonempty open set of analytic periodic tables admitting the
networks in Theorem~\ref{thm:network}, without reflection symmetry.
If the full-rank condition is omitted, a registered one-channel law can
remain unchanged while the periodic table changes by a nonisometric
lattice shear.
\end{proposition}
\begin{proof}
Start with a unit disk on the lattice $3\mathbb Z\times4\mathbb Z$.
Use the horizontal and vertical normal channels to its neighboring
copies, with labels $(1,0)$ and $(0,1)$. Their gaps are one and two.
The four contact occurrences form the registration on the single
obstacle vertex. The channels are isolated strict minima of the local
distance action, and their segments have positive clearance from all
other obstacles. These properties persist under sufficiently small
analytic $C^2$ perturbations of the support function and of the lattice:
the implicit-function theorem continues the normal contacts, and a
uniform diameter bound reduces the clearance check to finitely many
nearby translates. Arbitrary small nonsymmetric perturbations are
allowed. The two cycle labels retain rank two.

For the last assertion keep the unit disk fixed and use the lattice
basis $(3,0),(s,4)$, with $|s|$ sufficiently small, observing only the
horizontal channel. Its contact arcs, gap, local action and twist remain
unchanged. The covolume is always $12$, so the normalized phase area is
always $12-\pi$. Positive clearance gives a common small observation
collar. The two intrinsic local laws and their contact registration
are therefore exactly the same for all these tables. They do not depend
on the unseen second period. Nevertheless the lattices for $s=0$ and
small $s\ne0$ are not isometric: their shortest vector independent of
$(3,0)$ has length respectively $4$ and $\sqrt{16+s^2}$.
Thus local analytic continuation determines the obstacle in this example,
but not an unobserved lattice direction. Equation~\eqref{eq:latticerecover}
identifies exactly the additional information that removes this freedom.
\end{proof}
